% =============================================================================
% Daftar Pustaka
%
% PPTA IPB 2026: Harvard nama-tahun dengan gaya CSE edisi ke-9:
%   - In-text: (Author Year) atau (Author1 dan Author2 Year) atau
%              (Author1 et al. Year) untuk >2 penulis
%   - Daftar pustaka urut alfabet, hanging indent 1 cm, periode sebagai
%     pemisah field, tahun tanpa kurung.
%   - Maksimum lima penulis, lalu "et al." untuk sumber dengan >5 penulis.
%   - DOI dicantumkan bila tersedia; sumber daring memuat tanggal akses dan URL.
%
% Perintah yang bisa dipakai di teks (sintaks natbib kompatibel):
%   \citep{key}                  -> (Author Year)
%   \citep[5]{key}               -> (Author Year:5)
%   \citet{key}                  -> Author (Year)
%   \citep{key1,key2}            -> (Author1 Year1; Author2 Year2)
%   \citeauthor{key}             -> Author
%   \citeyear{key}               -> Year
%
% Backend: biber. Style didefinisikan di src/config/ipb-thesis.sty.
% Daftar pustaka di-deklarasikan di src/main.tex melalui \addbibresource.
% =============================================================================

\chapter*{DAFTAR PUSTAKA}
\addcontentsline{toc}{chapter}{DAFTAR PUSTAKA}

% Cetak daftar pustaka. Heading none agar tidak ada judul tambahan
% (judul "DAFTAR PUSTAKA" sudah ditangani oleh \chapter* di atas).
% SEMENTARA: \nocite mencetak semua pustaka kerja agar daftar bisa diperiksa.
% Hapus baris ini sebelum naskah akhir; hanya sumber yang diacu yang dicetak.
\nocite{Noviana2025,Pierce2025,Sahri2021,Sahri2023,Serov2026,Wilson2026,%
  Dares2025,Gonnerman2026,Smith2021,Fiedler2018,Correia2021,Bennington2020,%
  Pasanisi2024,Karimi2026,Sapag2025,Elith2010,Adriani2023,Bridges2023,%
  Beddoe2026,Huang2007,Saigusa2026,Patterson2025,Sequeira2012,Sequeira2014,%
  McKinney2012,Baez2020,PetatanRamirez2020,Milles2025,Miller2026,%
  Zhuang2021,Kouw2021,Zhang2024,Lalou2025,Leroy2018,MeyerPebesma2021,%
  Velazco2024,Moore2024,Liu2025,Schumacher2025,Matsui2026}
\printbibliography[heading=none]
